% !TEX root = ../../LAFROX_AnalisisFinanciero_2_20261113.tex
\section{Análisis comparativo de costos frente a precio}
\label{sec:e2-costos-precio}

El proyecto cuesta 194.487,71 UF y se ofrece a un precio neto de 228.763,92 UF, con un margen de 14,98\,\%. Las cifras de este capítulo son las del modelo financiero (LafroX, 2026f).

Los montos son netos y se expresan en UF, CLP y USD con los valores del Formulario E-24, una UF de \$40.000 y un dólar de \$900 (Distribuidora Puelche S.A., 2026a). Una UF equivale entonces a 44,44 dólares. Los montos en CLP y USD se redondean a la unidad. Las tablas mensuales, de sensibilidad y de innovaciones van solo en UF por su extensión y se convierten con los mismos valores.

\subsection{Cuadro maestro de costos}
\label{sec:e2-cuadro-costos}

Los costos de los meses 1 a 20 se cargan a la implementación, con 85.446,06 UF, y los de los meses 21 a 56 a la operación, con 109.041,65 UF. La Tabla~\ref{tab:e2-costos-maestro} los divide en costos directos de cada fase y costos indirectos.

% En esta tabla los filetes no admiten corte de página: solo se puede partir antes de un grupo.
{\renewcommand{\LFXfiletefila}{\noalign{{\color{lafroxLinea}\hrule height 0.3pt}\nobreak}}%
\begin{tablalafrox}{Cuadro maestro de costos por tipo y categoría, en UF, CLP y USD. Las sumas pueden diferir en la última cifra por redondeo. Fuente: modelo financiero, hojas COSTOS y RESUMEN.}{tab:e2-costos-maestro}{Y R{2.3cm} R{3.0cm} R{2.4cm}}
  {\cab{Categoría} & \cab{UF} & \cab{CLP} & \cab{USD}}
  \textbf{Costos directos de implementación} & & & \\
  Recursos humanos & 62.296,43 & 2.491.857.060 & 2.768.730 \\
  Capacitación, mentoría y certificación & 1.387,35 & 55.494.180 & 61.660 \\
  Infraestructura y equipamiento & 12.975,08 & 519.003.390 & 576.670 \\
  Servicios de proveedores & 4.472,14 & 178.885.729 & 198.762 \\
  Subtotal & 81.131,01 & 3.245.240.360 & 3.605.823 \\
  \noalign{\penalty0}\textbf{Costos directos de operación} & & & \\
  Recursos humanos & 74.273,19 & 2.970.927.600 & 3.301.031 \\
  Capacitación, mentoría y certificación & 1.933,53 & 77.341.020 & 85.934 \\
  Infraestructura y equipamiento & 671,66 & 26.866.393 & 29.852 \\
  Servicios de proveedores & 25.623,76 & 1.024.950.443 & 1.138.834 \\
  Subtotal & 102.502,14 & 4.100.085.456 & 4.555.651 \\
  \noalign{\penalty0}\textbf{Costos indirectos y overhead} & & & \\
  Gastos administrativos de implementación & 3.592,47 & 143.698.777 & 159.665 \\
  Gastos administrativos de operación & 6.225,30 & 249.012.090 & 276.680 \\
  Reservas de riesgo de implementación & 722,58 & 28.903.340 & 32.115 \\
  Reservas de riesgo de operación & 314,21 & 12.568.400 & 13.965 \\
  Subtotal & 10.854,57 & 434.182.607 & 482.425 \\
  \midrule
  \textbf{Total del proyecto} & \textbf{194.487,71} & \textbf{7.779.508.423} & \textbf{8.643.898} \\
\end{tablalafrox}}

Los recursos humanos son la mayor partida en las dos fases. Las 197.894,23 HH de los perfiles del T-15 cuestan 139.890,50 UF, y 3.320,88 UF de ese monto son horas de capacitación, mentoría y certificación, que se muestran aparte (LafroX, 2026c). El costo por hora de cada perfil está dentro del rango de su categoría en el Formulario E-26.

La operación cuesta más que la implementación porque mantiene la mesa de ayuda, el NOC y el SOC durante 36 meses y paga 25.623,76 UF en servicios de proveedores, entre nube, enlaces, el satelital de los camiones con frío, licencias y servicios de mapas y mensajería. En la implementación pesa la compra de infraestructura, 12.975,08 UF en equipos de red, seguridad y borde y en las obras de la sala técnica. Los equipos de terreno los compra el CLIENTE, quedan fuera del costo y del flujo, y su costo estimado figura en la hoja \texttt{HARDWARE DEL CLIENTE}.

Los costos indirectos suman 10.854,57 UF, de los cuales 9.817,77 UF son gastos administrativos y 1.036,79 UF reservas de riesgo. Las reservas pagan las horas de prevención, contingencia y gestión del plan de riesgos (LafroX, 2026e). La Tabla~\ref{tab:e2-gastos-adm} detalla los gastos administrativos.

\begin{tablalafrox}{Composición de los gastos administrativos en UF, CLP y USD. Fuente: modelo financiero, hoja GASTOS ADMINISTRATIVOS.}{tab:e2-gastos-adm}{Y R{2.3cm} R{3.0cm} R{2.4cm}}
  {\cab{Gasto} & \cab{UF} & \cab{CLP} & \cab{USD}}
  Overhead de la empresa & 6.048,00 & 241.920.000 & 268.800 \\
  Garantías & 1.422,20 & 56.887.863 & 63.209 \\
  Seguros & 2.179,33 & 87.173.129 & 96.859 \\
  Gastos notariales y de contratación & 5,30 & 212.000 & 236 \\
  Custodia de fuentes ante terceros & 162,95 & 6.517.875 & 7.242 \\
  \midrule
  \textbf{Total} & \textbf{9.817,77} & \textbf{392.710.867} & \textbf{436.345} \\
\end{tablalafrox}

De estos gastos, 3.592,47 UF caen en la implementación y 6.225,30 UF en la operación, según el mes de pago. El overhead es de 108 UF al mes durante todo el contrato y cubre el 25\,\% de la gerencia general, la administración y finanzas, la mitad del contador externo y el 30\,\% de la asesoría legal. Las garantías son las comisiones de las boletas de seriedad, fiel cumplimiento y correcto funcionamiento, y los seguros cubren responsabilidad civil, ciberriesgo y los equipos instalados en los sitios del CLIENTE (Distribuidora Puelche S.A., 2026a, arts.~35 a 38). Completan la tabla la firma del contrato en notaría y el depósito de los códigos fuente ante un tercero. Los intereses del crédito no están aquí; se registran en el financiamiento de la sección~\ref{sec:e2-curva-s}.

La hoja \texttt{RESUMEN} agrupa estas categorías en seis líneas (Tabla~\ref{tab:e2-costos-monedas}). Las líneas de recursos humanos incluyen las reservas de riesgo, y la de soporte suma también la capacitación de la operación. Las líneas de servicios incluyen los gastos administrativos, y la de soporte suma además la infraestructura de la operación.

\begin{tablalafrox}{Costo por línea del Resumen en UF, CLP y USD. Las sumas pueden diferir en la última cifra por redondeo. Fuente: modelo financiero, hoja RESUMEN.}{tab:e2-costos-monedas}{Y R{2.3cm} R{3.0cm} R{2.4cm}}
  {\cab{Línea} & \cab{UF} & \cab{CLP} & \cab{USD}}
  Implementación recursos humanos & 63.019,01 & 2.520.760.400 & 2.800.845 \\
  Implementación adquisiciones & 12.975,08 & 519.003.390 & 576.670 \\
  Implementación capacitaciones & 1.387,35 & 55.494.180 & 61.660 \\
  Implementación servicios & 8.064,61 & 322.584.506 & 358.427 \\
  \textbf{Total implementación} & \textbf{85.446,06} & \textbf{3.417.842.477} & \textbf{3.797.603} \\
  Soporte recursos humanos & 76.520,93 & 3.060.837.020 & 3.400.930 \\
  Soporte servicios & 32.520,72 & 1.300.828.926 & 1.445.365 \\
  \textbf{Total soporte} & \textbf{109.041,65} & \textbf{4.361.665.946} & \textbf{4.846.295} \\
  \midrule
  \textbf{Total del proyecto} & \textbf{194.487,71} & \textbf{7.779.508.423} & \textbf{8.643.898} \\
\end{tablalafrox}

\subsection{Precio de venta y márgenes por componente}
\label{sec:e2-precio-margenes}

Cada componente lleva su propio margen. La hora de cada perfil se cobra a 1,20 veces su costo, salvo en siete perfiles donde ese valor queda bajo la tarifa mínima del Formulario E-26 y se cobra la mínima (Distribuidora Puelche S.A., 2026a). Esos perfiles son el Líder de Datos, el Líder de Calidad y Pruebas, el Encargado de Seguridad de la Información y los ingenieros senior y junior de Datos y Bases y de DevOps y Plataforma. Las horas de capacitación se cobran con la tarifa del perfil que las dicta. Las adquisiciones, los servicios de proveedores y los gastos administrativos llevan un 10\,\% sobre el costo, ya que son compras a terceros y no trabajo del equipo. Las reservas de riesgo pasan al precio sin recargo. La Tabla~\ref{tab:e2-precio-monedas} muestra el precio de cada línea, en el mismo orden que la Tabla~\ref{tab:e2-costos-monedas}.

\begin{tablalafrox}{Precio de venta por línea del Resumen en UF, CLP y USD. Las sumas pueden diferir en la última cifra por redondeo. Fuente: modelo financiero, hojas INGRESOS y RESUMEN.}{tab:e2-precio-monedas}{Y R{2.3cm} R{3.0cm} R{2.4cm}}
  {\cab{Línea} & \cab{UF} & \cab{CLP} & \cab{USD}}
  Implementación recursos humanos & 76.155,21 & 3.046.208.364 & 3.384.676 \\
  Implementación adquisiciones & 14.272,59 & 570.903.729 & 634.337 \\
  Implementación capacitaciones & 1.664,83 & 66.593.016 & 73.992 \\
  Implementación servicios & 8.871,07 & 354.842.957 & 394.270 \\
  \textbf{Total implementación} & \textbf{100.963,70} & \textbf{4.038.548.066} & \textbf{4.487.276} \\
  Soporte recursos humanos & 92.027,42 & 3.681.096.800 & 4.090.108 \\
  Soporte servicios & 35.772,80 & 1.430.911.819 & 1.589.902 \\
  \textbf{Total soporte} & \textbf{127.800,22} & \textbf{5.112.008.619} & \textbf{5.680.010} \\
  \midrule
  \textbf{Precio neto del proyecto} & \textbf{228.763,92} & \textbf{9.150.556.685} & \textbf{10.167.285} \\
\end{tablalafrox}

El precio neto es de 228.763,92 UF, o \$9.150.556.685 y USD~10.167.285. La implementación se cobra en los doce hitos del Formulario E-25 y la operación en 36 cuotas de 3.550,01 UF (Distribuidora Puelche S.A., 2026a). La Tabla~\ref{tab:e2-margen} muestra el margen de cada fase.

\begin{tablalafrox}{Margen y rentabilidad del proyecto en UF, CLP y USD. Fuente: modelo financiero, hoja RESUMEN.}{tab:e2-margen}{Y R{2.3cm} R{3.0cm} R{2.4cm}}
  {\cab{Concepto} & \cab{UF} & \cab{CLP} & \cab{USD}}
  Margen de implementación & 15.517,64 & 620.705.590 & 689.673 \\
  Margen de soporte & 18.758,57 & 750.342.673 & 833.714 \\
  \textbf{Margen del proyecto} & \textbf{34.276,21} & \textbf{1.371.048.262} & \textbf{1.523.387} \\
  \midrule
  Rentabilidad sobre el precio & 14,98\,\% & 14,98\,\% & 14,98\,\% \\
\end{tablalafrox}

Las dos fases dejan margen, y el del proyecto llega a 34.276,21 UF. Como EBITDA sobre ingresos netos equivale a 14,98\,\%, bajo el tope de 20\,\% del Formulario E-24. Cada hito y cada cuota se facturan con un IVA de 19\,\%, desglosado en la Tabla~\ref{tab:e2-iva}.

\begin{tablalafrox}{Precio neto, IVA y precio total por fase en UF, CLP y USD. Fuente: modelo financiero, hojas RESUMEN y PARAMETROS.}{tab:e2-iva}{Y R{2.3cm} R{3.0cm} R{2.4cm}}
  {\cab{Concepto} & \cab{UF} & \cab{CLP} & \cab{USD}}
  Implementación, precio neto & 100.963,70 & 4.038.548.066 & 4.487.276 \\
  Implementación, IVA 19\,\% & 19.183,10 & 767.324.133 & 852.582 \\
  Operación, precio neto & 127.800,22 & 5.112.008.619 & 5.680.010 \\
  Operación, IVA 19\,\% & 24.282,04 & 971.281.638 & 1.079.202 \\
  \midrule
  \textbf{Precio neto del proyecto} & \textbf{228.763,92} & \textbf{9.150.556.685} & \textbf{10.167.285} \\
  \textbf{IVA del proyecto} & \textbf{43.465,14} & \textbf{1.738.605.770} & \textbf{1.931.784} \\
  \textbf{Precio total con IVA} & \textbf{272.229,06} & \textbf{10.889.162.455} & \textbf{12.099.069} \\
\end{tablalafrox}

La cuota de operación de 3.550,01 UF lleva 674,50 UF de IVA, y se factura en total 4.224,51 UF al mes. El flujo de caja y los indicadores usan los valores netos, ya que el IVA se traspasa al fisco.
